%% ma: L12-B06-0083
%% lop: 12
%% chuong: 2
%% bai: 6
%% dang: 12.6.22
%% loai: TN4
%% muc: VDC
%% dap_an: "A"
%% nguon: "CĐ2.6.pdf giáo viên gửi 27/10/2026 (M07, Lê Minh Khôi – Nguyễn Anh Khôi, Chương II Chủ đề 6 Ứng dụng vào các bài toán thực tế), Câu 26"
%% kien_thuc: "Bài 6 (tổng vectơ, bình phương vô hướng), lớp 10 (bất đẳng thức, tam thức bậc hai)"
%% trang_thai: chinh-thuc
%% ghi_chu: "Giáo viên đã duyệt 27/10/2026. Đã bổ sung giả thiết độ lớn lực F2 thay đổi (đề gốc chưa nêu). Đề gốc dùng chung giả thiết cho Câu 25–26; đã chép vào từng câu; bổ sung 'Khi độ lớn của lực F2 thay đổi' cho rõ đại lượng biến thiên"
\begin{ex}
Trong không gian, cho bốn sợi dây buộc với nhau tại điểm $I$. Điểm $I$ cân bằng dưới tác dụng của bốn lực dọc theo bốn sợi dây lần lượt là $\vec F_1$, $\vec F_2$, $\vec F_3$, $\vec F_4$. Biết rằng $|\vec F_1|=180$ N, $|\vec F_3|=220$ N, góc tạo bởi hai lực $\left(\vec F_2,\vec F_4\right)=120^\circ$, góc tạo bởi hai lực $\left(\vec F_1,\vec F_3\right)=50^\circ$. Đơn vị lực trong bài toán tính theo Newton. Khi độ lớn của lực $\vec F_2$ thay đổi, để hệ thống duy trì trạng thái cân bằng với các góc đã cho, độ lớn của lực $\vec F_4$ có thể đạt giá trị lớn nhất bằng bao nhiêu (làm tròn kết quả đến hàng đơn vị)?
\choice{\True $419$ N}{$726$ N}{$363$ N}{$209$ N}
\loigiai{$R=\left|\vec F_1+\vec F_3\right|\approx362{,}9$ N cố định. Với $a=|\vec F_2|$, $b=|\vec F_4|$: $a^2-ab+b^2=R^2$.

$R^2=\left(a-\dfrac b2\right)^2+\dfrac34b^2\ge\dfrac34b^2\Rightarrow b\le\dfrac{2R}{\sqrt3}\approx419$ N (đạt khi $a=\dfrac b2$).}
\end{ex}
